% !TeX program = xelatex
% UTF-8; standalone: xelatex "common cold.tex" from this directory.
% Inclusion: define \CommonColdInNotebook before \input in a XePersian document.
% Source: Google Drive Notebook/common cold; nine images in filename order.
% pic01.jpg = 20261012_173955944.jpg; 1EuLaIWEFL7iOx2RhZPUW5oOkgbbKUa7Z
% pic02.jpg = 20261012_174001039.jpg; 1LXBwEXyLCFP0Tiqlj8tAUymfwZ1mEMD7
% pic03.jpg = 20261012_174005344.jpg; 1Q1qJ1QawMZGe3AETDQ_t8uRnudO8jEKw
% pic04.jpg = 20261012_174010011.jpg; 1CGEInzK8lhMCktl-w6hicWytqybPX07x
% pic05.jpg = 20261012_174017688.jpg; 1H9jZotigCOYmAin45zhMuwxLZofZvLub
% pic06.jpg = 20261012_174022537.jpg; 1d1uQXd8qq7-FQ_xa3lppoPkq9ij5gsm7
% pic07.jpg = 20261012_174027760.jpg; 1Hw5_DkFUFjcXVMNb3ajPA5fi7bXf4UGF
% pic08.jpg = 20261012_174033748.jpg; 1eRIYgewe7Wh_7JHtK2scMBDXep1fQBwB
% pic09.jpg = 20261012_174038826.jpg; 1oheOIINORFmKWgxboNWlfXQIXVxxTX77
% Images intentionally use the user-requested spelling: ped-infectioud diseases.
\ifdefined\CommonColdInNotebook
  \let\CCfinish\relax
\else
  \def\CCfinish{\end{document}}
  \documentclass[12pt,a4paper]{article}
  \usepackage[margin=22mm]{geometry}
  \usepackage{amsmath,amssymb}
  \usepackage{graphicx}
  \usepackage{xepersian}
  \IfFontExistsTF{Amiri}{\settextfont{Amiri}}{\settextfont{DejaVu Sans}}
  \setlatintextfont{DejaVu Serif}
  \setlength{\parindent}{0pt}
  \setlength{\parskip}{5pt}
  \setlength{\emergencystretch}{3em}
  \begin{document}
\fi
\providecommand{\CCunclear}[1]{\textbf{[خوانش نامطمئن: #1]}}
\providecommand{\CCpage}[2]{\clearpage\section*{تصویر #1}\lr{\texttt{#2}}\par
  \textbf{سربرگ و حاشیه:} ۴۲۸؛ \lr{Subject:} (خالی)؛ \lr{Date: / /}؛ شمارهٔ صفحه: #1؛ نشان پایین برگه: \lr{YASHA}.\par}
\providecommand{\CCrange}[2]{\lr{#1--#2}}
\begin{center}
{\Large\bfseries \lr{The Common Cold}}\par
\lr{Nelson 22nd}\par
۴۲۸
\end{center}
\textbf{یادداشت رونویسی:} متن و اختصارات منبع، به ترتیب نه تصویر، حفظ شده‌اند. فلش‌ها به صورت فهرست و رابطه بازنویسی شده‌اند. گزاره‌های پزشکی منبع اصلاح نشده‌اند. خوانش نامطمئن با برچسب مشخص شده است. تصاویر اصلی کامل در پیوست آمده‌اند.

\clearpage\section*{تصویر ۱}\lr{\texttt{pic01.jpg}}\par
\textbf{سربرگ و حاشیه:} \lr{The Common Cold}، ۴۲۸؛ \lr{Subject: Nelson 22nd}؛ \lr{Date: / /}؛ شمارهٔ صفحه: ۱؛ \lr{YASHA}.
\begin{itemize}
\item علائم بالینی اصلی: رینوره، انسداد بینی.
\item در \lr{Common Cold} علائم سیستمیک یا وجود ندارد یا خفیف می‌باشد: سردرد، میالژی، تب.
\item نام دیگر: رینوسینوزیت.
\end{itemize}
\subsection*{اتیولوژی}
\begin{itemize}
\item بیش از ۲۰۰ نوع رینوویروس انسانی (\lr{HRV}) (عامل ۵۰\% موارد).
\item \lr{RSV}، \lr{MPV} (متاپنوموویروس)، \lr{PIV} (پاراآنفلوانزا)، کروناویروس، آدنوویروس.
\end{itemize}
\subsection*{اپیدمیولوژی}
\begin{itemize}
\item در تمام طول سال ممکن است ولی بیشتر در ابتدای پاییز تا اواخر بهار.
\item کودکان کم‌سن معمولاً بین \CCrange{6}{8} نوبت در سال سرما می‌خورند.
\item \CCrange{10}{15}\% کودکان بیش از ۱۲ نوبت در سال.
\item \CCrange{2}{3} نوبت در سال در بزرگسالان.
\item بروز موارد سرماخوردگی در ۳ سال اولی که کودک به مهدکودک / مدرسه می‌رود بالاتر از حالتی است که در خانه بماند.
\item وقتی به مدرسه می‌روند آنهایی که مهدکودک رفته‌اند کمتر مریض می‌شوند.
\end{itemize}
کمبود لکتین متصل‌شونده به مانوز، با اختلال در ایمنی ذاتی، با افزایش موارد ابتلا به سرماخوردگی در یک سال همراه است.
\begin{latin}\noindent Mannose-binding lectin deficiency\end{latin}

\CCpage{۲}{pic02.jpg}
\subsection*{پاتوژنز سرماخوردگی}
\begin{itemize}
\item ویروس‌های عامل سرماخوردگی از ۳ مسیر منتشر می‌شوند:
\begin{enumerate}
\item تماس مستقیم دست.
\item استنشاق ویروس \lr{airborne}.
\item ذرات ویروس بزرگ‌تر مثل عطسه.
\end{enumerate}
\item سرماخوردگی با \lr{HRV} و آدنوویروس: ایمنی سروتایپ اختصاصی ایجاد می‌کنند.
\item آنفلوانزا آنتی‌ژن‌های سطحی‌اش را تغییر می‌دهد.
\item کروناویروس تایپ‌های زیادی دارد که در برابر این ایمنی ایجاد می‌کند.
\item حتی وقتی ایمنی علیه ویروس شکل گرفته باز هم فرد سرماخورده با آن ویروس ممکن است، چرا که عمر \lr{IgA} موجود در مخاط کوتاه است و ویروس به صورت موضعی در آنجا تکثیر می‌شود بدون آنکه ایمنی را تحریک کند.
\item ولی شدت بیماری خفیف‌تر است.
\item ویروس آنفلوانزا و آدنوویروس اپی‌تلیوم را تخریب می‌کنند و آسیب پاتولوژیک ایجاد می‌نمایند؛ در مقابل \lr{HRV}، کروناویروس، \lr{RSV} آسیب ایجاد نمی‌نمایند.
\item در مواجهه با ویروس واکنش التهابی حاد رخ می‌دهد.
\item \lr{IL-8}: جذب \lr{PMN}؛ عامل علائم بالینی سرماخوردگی.
\item خود ویروس عامل علائم بالینی نیست.
\item انتشار ویروس \CCrange{3}{5} روز پس از ورود ویروس به بدن ادامه دارد.
\item انتشار خفیف ممکن است تا ۳ هفته ادامه یابد، حتی اگر بیمار کامل خوب شده باشد.
\item زمانی که علائم بالینی شروع می‌شود معمولاً \CCunclear{شدینگ} خاتمه می‌یابد.
\item التهاب ممکن است دهانهٔ سینوس‌ها و شیپور استاش را ببندد: مستعد سینوزیت باکتریال و اوتیت مدیا.
\end{itemize}

\CCpage{۳}{pic03.jpg}
\subsection*{دیگر یافته‌های بالینی در سرماخوردگی}
\begin{itemize}
\item در \lr{HRV} و کروناویروس: ویزینگ / برونشیولیت.
\item آنفلوانزا و پاراآنفلوانزا: \lr{Croup}.
\item آدنوویروس: تب فارنگوکنژنکتیوال (کنژنکتیویت، التهاب بینی، التهاب فارنکس).
\item انتروویروس‌ها: هرپانژین (تب و پاپول‌های اریتماتو روی اوروفارنکس).
\end{itemize}
\subsection*{تظاهرات بالینی در سرماخوردگی}
\begin{itemize}
\item در شیرخواران: تب و دیس‌چارژ بینی.
\item تب در کودکان بزرگ‌تر نادر است.
\item علائم سرماخوردگی \CCrange{1}{3} روز پس از مواجهه با ویروس شروع می‌شود.
\item اولین علامت معمولاً سوزش و خارش گلو است؛ در ادامه انسداد بینی و آبریزش.
\item گلودرد معمولاً به سرعت رفع می‌شود و در روز دوم و سوم بیماری علائم بینی غالب می‌شوند.
\item سرفه در \lr{\(2/3\)} موارد سرماخوردگی وجود دارد در کودکان و معمولاً پس از علائم بینی آغاز می‌شود.
\item سرفه ممکن است \CCrange{1}{3} هفته پس از رفع سایر علائم نیز ادامه یابد.
\item \lr{HRV} و کروناویروس معمولاً با تب و علائم سیستمیک همراه خیلی کمتری دارند.
\item استفراغ و اسهال ناشایع است.
\item معمولاً یک هفته طول می‌کشد.
\item حدود ۱۰\% موارد برای دو هفته طول می‌کشند.
\item تغییر رنگ و قوام مایع ترشح بینی در سیر بیماری رخ می‌دهد.
\item ارتباطی با عفونت باکتریال یا سینوزیت ندارد.
\item مطرح‌کنندهٔ تجمع \lr{PMN}.
\end{itemize}

\CCpage{۴}{pic04.jpg}
\begin{itemize}
\item فشار غیرطبیعی گوش میانی در طی سرماخوردگی شایع است.
\item لنفادنوپاتی بزرگ در قدام سرویکال ممکن است دیده شود.
\item کنژنکتیویت نیز ممکن است رخ دهد.
\end{itemize}
\subsection*{تشخیص سرماخوردگی}
\begin{itemize}
\item مهم‌ترین اقدام \lr{R/O} سایر \lr{DDX}های سرماخوردگی است.
\item رینیت آلرژیک: عطسه و خارش پرزیستنت هستند؛ ائوزینوفیل بینی، \lr{Hansel stain}.
\item رینیت وازوموتور: به دنبال محرک‌ها، تغییرات دما، تغییرات آب و هوا.
\item رینیت مدیکامنتوزا: شرح حال مصرف وازوکانستریکتان بینی.
\item جسم خارجی بینی: یک‌طرفه، ترشحات بدبو (\lr{foul smelling}).
\item سینوزیت: تب عودشونده، درد صورت و ادم پری‌اوربیت و رینوره ادامه‌دار، سرفه بیش از ۱۴ روز.
\item استرپتوکوکوزیس: ترشحات نازال موکوپورولنت که \CCunclear{پوست بینی را خراش می‌دهد}، رینوره ندارد.
\item پرتوزیس: شروع سرفه پاروکسیسمال شدید.
\item سیفلیس مادرزادی: رینورهٔ ماندگار با شروع در ۳ ماه اول زندگی.
\end{itemize}
\subsection*{یافته‌های آزمایشگاهی در سرماخوردگی}
\begin{itemize}
\item آزمایش روتین در تشخیص و درمان کمک‌کننده نیست.
\item برای \lr{R/O} رینیت آلرژیک: \lr{Nasal Smear} از نظر \lr{Hansel stain}.
\item در اسمیر نازال \lr{PMN} دیده شده است: شاخصی از سرماخوردگی بودن عارضه است و مطرح‌کنندهٔ عفونت باکتریال نیست.
\end{itemize}

\CCpage{۵}{pic05.jpg}
تصویربرداری اندیکاسیون ندارد.
\begin{itemize}
\item اگر انجام شود احتمالاً ابنورمالیتی‌هایی دیده می‌شوند در سینوس‌های پارانازال.
\end{itemize}
شناسایی پاتوژن‌ها (ویروس) با \lr{PCR}، کشت، تشخیص آنتی‌ژن و سرولوژی ممکن است.
\begin{itemize}
\item در سرماخوردگی به دنبال شناسایی نوع ویروس نیستیم.
\item مگر زمانی که بخواهیم درمان ضدویروسی شروع کنیم، مثل آنفلوانزا.
\end{itemize}
کشت مثبت باکتریال از نمونه‌های نازوفارنژیال به معنی عفونت باکتریال نیست و مطرح‌کنندهٔ علت سینوزیت هم نمی‌باشد.
\subsection*{درمان سرماخوردگی}
حمایتی و انتظاری: \lr{Supportive and Anticipatory}.
\subsection*{درمان ضدویروسی}
\begin{itemize}
\item آنفلوانزا: مهارکنندهٔ نورآمینیداز: اوسلتامیویر، زانامیویر، پرامیویر.
\item مهارکنندهٔ پلی‌مراز اسید اندونوکلئاز: بالوکساویر ماربوکسیل.
\item درمان باید ظرف \lr{48h} از شروع علائم آغاز شود.
\end{itemize}
درمان آنتی‌باکتریال هیچ فایده‌ای در سرماخوردگی ندارد: نباید داده شود.
\subsection*{علامت‌درمانی}
\begin{itemize}
\item مصرف مایعات کافی برای جلوگیری از دهیدره شدن و هومیدیفایر برای رقیق شدن ترشحات بینی و مجاری تنفسی.
\item مایعات گرم گلو را نرم و روان می‌کنند.
\end{itemize}

\CCpage{۶}{pic06.jpg}
\begin{itemize}
\item \lr{Topical Nasal Saline}: ترشحات را می‌برد.
\item ایریگیشن با سالین هم علائم را کاهش می‌دهد.
\item هوای مرطوب گرم: ترشحات را شل می‌کند؛ توصیه نشده است.
\item استفاده از داروهای \lr{OTC} در کودکان محل بحث است.
\item اگرچه در بزرگسالان مؤثر هستند، ولی در کودکان مطالعاتی نداریم.
\item تصمیم به استفاده باید بر اساس سود و زیان باشد.
\item قرص‌های ترکیبی ندهیم بلکه بر اساس علامت اصلی بدهیم.
\item قرص مکیدنی \lr{Zinc (lozenges)} در کودک سالم اگر ظرف ۲۴ ساعت اولیهٔ شروع علائم آغاز شود می‌تواند طول مدت بیماری را کاهش دهد ولی شدت آن را نه.
\item عوارض شامل کاهش چشایی، مزهٔ بد، تهوع.
\end{itemize}
تب معمولاً در سرماخوردگی وجود ندارد پس دیگر ضدتب اندیکاسیون ندارد.

در میالژی و سردرد سرماخوردگی می‌توان از \lr{NSAID} استفاده کرد.
\subsection*{انسداد بینی در سرماخوردگی}
\begin{itemize}
\item آدرنرژیک خوراکی یا موضعی به عنوان ضداحتقان.
\item موضعی: اکسی‌متازولین، زایلو‌متازولین، نفازولین؛ قطره یا اسپری بینی.
\item زیر ۶ سال توصیه نمی‌شوند.
\item جذب سیستمیک ایمیدازولین‌ها: عارضهٔ مشابهٔ داروی کاردیوواسکولار، هایپوتنشن، کوما، کولاپس دارد.
\item در صورت استفادهٔ طولانی‌مدت ریباند دارد: رینیت مدیکامنتوزا.
\end{itemize}

\CCpage{۷}{pic07.jpg}
\begin{itemize}
\item آدرنرژیک‌های خوراکی اثر کمتری دارند؛ عارضه: تحریک سیستمیک، \lr{CNS}.
\item سودوافدرین مؤثرترین ضداحتقان خوراکی است.
\item قطرهٔ بینی (سالین) در هر سن مجاز است و مؤثر می‌باشد.
\end{itemize}
\subsection*{رینوره در سرماخوردگی}
\begin{itemize}
\item آنتی‌هیستامین نسل \lr{I} رینوره را به میزان \CCrange{25}{30}\% کاهش می‌دهد.
\item اثر آنها به علت اثر آنتی‌کولینرژیک‌شان است، نه آنتی‌هیستامین آنها؛ به همین علت نسل \lr{II} مؤثر نیستند.
\item ایپراتروپیوم بروماید: آنتی‌کولینرژیک موضعی.
\item اثر مشابهٔ نسل \lr{I} ولی خواب‌آور نیست.
\item عارضهٔ شایع: تحریک بینی و خونریزی.
\end{itemize}
\subsection*{گلودرد در سرماخوردگی}
\begin{itemize}
\item معمولاً خیلی جدی نیست و نیاز به درمان ندارد.
\item استامینوفن طی عفونت \lr{HRV} منجر به سرکوب پاسخ خنثی‌کنندهٔ \lr{Ab} می‌شود که البته از نظر بالینی بی‌اهمیت است.
\item آسپرین نباید به کودک با عفونت تنفسی داده شود: سندرم \lr{Reye} در کودک با آنفلوانزا.
\end{itemize}
\subsection*{سرفه در سرماخوردگی}
\begin{itemize}
\item معمولاً \CCunclear{ضددرمانی} ندارد.
\item ممکن است در زمینهٔ \lr{PND} باشد در برخی از بیماران.
\item در این بیماران سرفه در زمان شدیدترین علائم بینی، در شدیدترین فرم خودش است؛ این گروه از آنتی‌هیستامین نسل \lr{I} سود می‌برند.
\end{itemize}

\CCpage{۸}{pic08.jpg}
\begin{itemize}
\item آبنبات مکیدنی ضدسرفه ممکن است مؤثر باشد.
\item در کودکان بی‌خطر است: در کودکانی که ریسک آسپیره کردن ندارند (بالای ۶ سال).
\item عسل اثر متوسطی بر کنترل سرفه‌های شبانه دارد.
\item در بالای ۱ سال ایمن است.
\end{itemize}
اگر سرفه طول بکشد (بیش از چند هفته):
\begin{latin}\noindent Virus induced reactive airway disease\end{latin}
\begin{itemize}
\item درمان با برونکودیلاتور.
\end{itemize}
دکسترومتورفان هیچ اثری بر سرفه در سرماخوردگی ندارد.

اکسپکتورانت‌ها مثل \lr{guaifenesin} اثری بر سرفه ندارند.

ویتامین \lr{C}، بخور گرم: بی‌اثر. اکیناسه باید \CCunclear{با قید استفاده باشد}.

\lr{Persistent acute purulent rhinitis} زیر ۱۰ روز در هیچ سنی نیاز به آنتی‌بیوتیک ندارد.
\subsection*{عوارض سرماخوردگی}
\begin{itemize}
\item \lr{AOM} شایع‌ترین عارضه است.
\item پس از چند روز سرماخوردگی: تب، درد گوش.
\item \CCrange{5}{30}\% بچه‌ها با سرماخوردگی.
\item در شیرخواران شایع‌تر است.
\item درمان علامتی سرماخوردگی تأثیری بر رخ دادن یا ندادن \lr{AOM} ندارد.
\end{itemize}

\CCpage{۹}{pic09.jpg}
\subsection*{سینوزیت باکتریال حاد}
\begin{itemize}
\item \CCrange{5}{9}\% کودکان و \lr{0.5--2}\% بزرگسالان مبتلا به سرماخوردگی. واژهٔ «کودکان» در بخش دوم خط خورده و با «بزرگسالان» جایگزین شده است.
\item رینوره یا سرفهٔ روزانهٔ ماندگار بدون تغییر برای \CCrange{10}{14} روز.
\item درمان سرماخوردگی تأثیری بر رخ دادن یا ندادن \lr{ABS} ندارد.
\end{itemize}
\subsection*{اگزاسربیشن آسم}
\begin{itemize}
\item یک عارضهٔ جدی سرماخوردگی است.
\end{itemize}
پروفیلاکسی دارویی یا ایمونولوژیک برای سرماخوردگی وجود ندارد.
\begin{itemize}
\item \lr{Palivizumab}: جلوگیری از عفونت \lr{RSV} در شیرخواران \lr{high risk}.
\item تأثیری بر \lr{URI} ناشی از \lr{RSV} ندارد.
\item ویتامین \lr{C}، سیر، اکیناسه: بی‌اثر بر پیشگیری از سرماخوردگی.
\end{itemize}

\clearpage\section*{پیوست: تصاویر کامل منبع}
% Relative paths when compiling inside the topic directory; repository-root fallback.
\ifdefined\CommonColdInNotebook
  \def\CCimagepath{pediatrics/ped-infectioud diseases/common cold/images/}
\else
  \IfFileExists{../../ped-infectioud diseases/common cold/images/pic01.jpg}
    {\def\CCimagepath{../../ped-infectioud diseases/common cold/images/}}
    {\def\CCimagepath{pediatrics/ped-infectioud diseases/common cold/images/}}
\fi
\newcommand{\CCoriginal}[1]{\clearpage\begin{center}
  \textbf{تصویر اصلی \lr{#1}}\par\smallskip
  \includegraphics[width=\textwidth,height=0.89\textheight,keepaspectratio]{\CCimagepath pic#1.jpg}
\end{center}}
\CCoriginal{01}
\CCoriginal{02}
\CCoriginal{03}
\CCoriginal{04}
\CCoriginal{05}
\CCoriginal{06}
\CCoriginal{07}
\CCoriginal{08}
\CCoriginal{09}
\CCfinish
